%% 402Pilot course project report.
%% Six-page self-contained version; technical_report.tex keeps proofs,
%% sensitivity sweeps, detailed ablations, and the x402 witness.

\documentclass[sigconf, nonacm]{acmart}
\raggedbottom
\settopmatter{printacmref=false, printccs=false, printfolios=true}
\renewcommand\footnotetextcopyrightpermission[1]{}

\usepackage{booktabs}
\usepackage{tabularx}
\usepackage{amsmath}
\usepackage{graphicx}
\usepackage{algorithm}
\usepackage[noend]{algpseudocode}
\usepackage{balance}
\usepackage{xurl}

\newcommand{\ut}{u_t}
\newcommand{\ctilde}{\tilde c_t}
\newcommand{\rt}{r_t}
\newcommand{\lamn}{\lambda_{\mathrm{norm}}}
\DeclareRobustCommand{\padct}{PA\nobreakdash-DCT}
\newcommand{\thetaq}[2]{\theta_q^{#1,#2}}
\newcommand{\thetac}[2]{\theta_c^{#1,#2}}
\newcommand{\Sigmaq}[2]{\Sigma_q^{#1,#2}}
\newcommand{\Sigmac}[2]{\Sigma_c^{#1,#2}}
\newcommand{\hatu}[2]{\hat u^{#1,#2}}
\newcommand{\hatc}[2]{\hat c^{#1,#2}}
\newcommand{\bdev}{\mathrm{burn\_dev}}

\AtBeginDocument{%
  \providecommand\BibTeX{{Bib\TeX}}}

\setcopyright{none}
\copyrightyear{2026}
\acmYear{2026}
\acmDOI{}
\acmConference[Course Project Report]{Course Project Report}{2026}{}
\acmISBN{}

\begin{document}

\title[402Pilot Course Report]{\texorpdfstring{402Pilot: An x402 Decision Layer for\\Autonomous Agent Micropayments}{402Pilot: An x402 Decision Layer for Autonomous Agent Micropayments}}
\subtitle{Course Project Report}

\author{Yin Li}
\authornote{Student ID: 20489800.}
\email{yligt@connect.hkust-gz.edu.cn}

\author{Yanbo He}
\authornote{Student ID: 50040464.}
\email{yhe376@connect.hkust-gz.edu.cn}

\author{Wanting Xiong}
\authornote{Student ID: 50022371.}
\email{wxiong945@connect.hkust-gz.edu.cn}
\affiliation{%
  \institution{The Hong Kong University of Science and Technology (Guangzhou)}
  \city{Guangzhou}
  \state{Guangdong}
  \country{China}
}
\renewcommand{\shortauthors}{Li, He, and Xiong}

\begin{abstract}
Programmable-payment protocols such as x402 make per-request agent payments executable, but they do not decide which service a budgeted agent should buy. We study this missing buyer-side layer as \emph{agent-native payment decision-making}: contextual provider selection under a draining wallet, uncertain service outcomes, and non-stationary reliability or prices. We propose \emph{402Pilot}, a protocol-agnostic decision layer above payment execution, and instantiate it with \emph{\padct}, a payment-aware discounted contextual Thompson-sampling policy. \padct{} maintains separate discount-tracked posteriors over service utility and realized cost, allowing both reliability drift and price drift to be tracked from post-payment feedback. We evaluate on \emph{402Pilot-Bench}, a frozen-replay benchmark with 823 tasks, five provider pipelines, and three market scenarios. \padct{} pays an exploration cost in the stationary setting, but provides the most robust adaptive trade-off under non-stationarity: during a mid-tier outage it improves full-horizon mean quality by 15 percentage points over the cheapest fixed provider while remaining budget-feasible, and under a premium price promotion it improves over the reliable mid-tier baseline on quality, ROI, and payment-aware gap.
\end{abstract}

\keywords{Agent payments, x402, micropayments, contextual bandits, non-stationary learning, Thompson sampling}

\maketitle

\section{Introduction}
\label{sec:intro}

x402-style payment protocols expose machine-readable payment requirements, per-request settlement, and increasingly dynamic service metadata, including pricing information~\citep{coinbase2025x402v2}. These capabilities make agent micropayments executable, but execution alone does not answer the buyer's decision question: how should an autonomous agent with a finite wallet decide which service is worth buying in a changing market?

Several lines of prior work supply pieces of this problem, but not the full buyer-side loop. Budgeted bandits model resource-constrained learning~\citep{xia2015bts,badanidiyuru2018bandits}; discounted methods address non-stationarity~\citep{garivier2011discounted,qi2023dsts}; contextual Thompson sampling estimates task-conditioned value~\citep{agrawal2013contextualts}; and LLM routing studies quality--cost tradeoffs across heterogeneous models~\citep{chen2023frugalgpt,ong2024routellm,ding2024hybridllm,mixllm2024,llmbandit2025}. These ingredients are usually analyzed outside a protocol-mediated quote--pay--observe loop in which the buyer sees payable endpoints, pays one of them, and updates from realized prices and service outcomes. Payment protocols expose the execution substrate, and auction mechanisms study strategic allocation~\citep{patra2026reverseauction}; neither provides an online spending policy for a single autonomous buyer.

We call this setting \emph{agent-native payment decision-making}. It differs from human-operated provider selection in three ways. First, the policy is part of an autonomous system, so routine operation cannot rely on a human resetting state or curating providers after each change. Second, the wallet drains with every paid call; each payment is irreversible and reveals only the chosen provider's outcome. Third, market conditions change: service quality, reliability, and effective price can shift while the policy observes them only through post-payment feedback. These properties make autonomous micropayment an online economic decision problem, not only a settlement problem.

This report proposes \emph{402Pilot}, a narrow decision layer that separates provider selection from payment execution. The payment stack handles how to pay; 402Pilot decides which affordable service is worth buying and updates that decision rule from post-payment feedback. We instantiate the layer with \emph{\padct} (Payment-Aware Discounted Contextual Thompson Sampling), a learning policy that maintains two discount-tracked beliefs for each provider--task bucket: one over service utility and one over realized cost. Its ranking rule combines wallet pressure, task context, discounted updating, and posterior sampling, allowing the same policy to react to reliability drift and price drift without hard-coding provider labels or assuming fixed prices.

We make three contributions. \textbf{C1}: we formulate agent-native payment decision-making and introduce 402Pilot as a protocol-agnostic buyer-side decision layer. \textbf{C2}: we instantiate the layer with \padct, a payment-aware discounted contextual Thompson-sampling policy with separate utility and cost beliefs. \textbf{C3}: we construct 402Pilot-Bench, a frozen-replay benchmark that evaluates buyer policies under stationary operation, reliability shocks, and price shocks.

\section{Related Work}
\label{sec:related-work}

\paragraph{Programmable payment protocols.}
x402 embeds payment into HTTP and exposes payment requirements to software agents; its V2 direction adds discovery, dynamic recipients, extensible schemes, and service metadata~\citep{coinbase2025x402v2}. AP2 provides payment-agnostic authorization for agent-led commerce~\citep{google2025ap2}; A2A--x402 and A402 bind agent-to-agent service execution to cryptocurrency-style settlement~\citep{a2ax402,a402}; and recent surveys organize agent payments around discovery, authorization, execution, and accounting~\citep{sok2026agentpayments}. 402Pilot is complementary: it treats these systems as rails and studies the buyer's spending policy above them.

\paragraph{Bandit ingredients and LLM routing.}
Budgeted Thompson Sampling and bandits with knapsacks model resource constraints~\citep{xia2015bts,badanidiyuru2018bandits}; contextual bandits condition value on task features~\citep{agrawal2013contextualts,li2010linucb}; and discounted or sliding-window methods track non-stationary feedback~\citep{garivier2011discounted,trovo2020sliding,qi2023dsts}. Cost-aware LLM routing systems such as FrugalGPT, RouteLLM, Hybrid LLM, MixLLM, and LLM Bandit also study quality--cost tradeoffs~\citep{chen2023frugalgpt,ong2024routellm,ding2024hybridllm,mixllm2024,llmbandit2025}. The key distinction is the information contract: routing systems usually assume published prices or offline labels, while 402Pilot learns from receipt-revealed charges and chosen-arm outcomes under a draining wallet.

\begin{table}[t]
\centering
\small
\setlength{\tabcolsep}{5pt}
\caption{Bandit ingredients related to PA-DCT. Budget denotes resource constraints, Drift non-stationary adaptation, and Context task-conditioned value estimation.}
\label{tab:bandit-components-short}
\begin{tabularx}{\columnwidth}{@{}>{\raggedright\arraybackslash}Xccc@{}}
\toprule
Method & Budget & Drift & Context \\
\midrule
BTS / BwK~\citep{xia2015bts,badanidiyuru2018bandits} & \checkmark & -- & -- \\
Linear CBwK~\citep{agrawal2016linearcbwk} & \checkmark & -- & \checkmark \\
DS-TS / Discounted UCB~\citep{qi2023dsts,garivier2011discounted} & -- & \checkmark & -- \\
Contextual TS / LinUCB~\citep{agrawal2013contextualts,li2010linucb} & -- & -- & \checkmark \\
\textbf{\padct{} (ours)} & \checkmark & \checkmark & \checkmark \\
\bottomrule
\end{tabularx}
\end{table}

\section{Problem and Layer}
\label{sec:problem-layer}

We model repeated provider selection by a single autonomous buyer with a finite wallet. There are $K$ payable providers and a horizon $T$. At round $t$, a task context $\mathbf{x}_t$ arrives, the buyer has remaining budget $B_t$, and provider $a$ has effective charge $p_{a,t}>0$. The wallet/payment interface returns the affordable set
\begin{equation}
\label{eq:affordable-set-short}
\mathcal{A}_t=\{a\in\{1,\ldots,K\}:p_{a,t}\le B_t\}.
\end{equation}
If $\mathcal{A}_t$ is empty, the run stops. Otherwise the policy chooses $a_t\in\mathcal{A}_t$, pays for that provider, and receives one feedback sample $(q_t,f_t,c_t)$, where $q_t\in[0,1]$ is service quality, $f_t\in\{0,1\}$ marks non-delivery, and $c_t=p_{a_t,t}$ is the realized charge. The wallet evolves as
\begin{equation}
\label{eq:wallet-short}
B_{t+1}=B_t-c_t .
\end{equation}
The policy does not observe unchosen outcomes, and the decision-layer contract only assumes access to the affordable set before payment; selected calls reveal numeric charges through receipts.

Service delivery is summarized as utility $\ut=q_t-\nu f_t$. We combine utility with normalized realized cost and wallet pressure:
\begin{equation}
\label{eq:reward-short}
\begin{aligned}
\ctilde&=\mathrm{clip}(c_t/c_{\max},0,1),\\
\lamn&=\lambda_t/(1+\lambda_t),\\
\rt&=(1-\lamn)\ut-\lamn\ctilde .
\end{aligned}
\end{equation}
For $0\le\nu\le1$, $\rt\in[-1,1]$. The pressure multiplier $\lambda_t$ grows when spending is ahead of the linear plan:
\begin{equation}
\label{eq:lambda-short}
\begin{aligned}
\mathrm{burn}_t&=\frac{(B-B_t)/B}{t},\\
\bdev_t&=\frac{\mathrm{burn}_t-1/T}{1/T},\\
\lambda_t&=\lambda_0\exp(\alpha\bdev_t).
\end{aligned}
\end{equation}
Because $\lambda_t$ depends on the policy's own wallet trajectory, evaluation reports payment-aware gap together with quality, budget use, and ROI rather than treating raw cumulative $\rt$ as a single leaderboard score.

\begin{figure*}[t]
\centering
\includegraphics[width=\linewidth,trim=0 55bp 0 40bp,clip]{figures/402pilotlayer.png}
\caption{The 402Pilot per-round decision flow. A task and quoted candidates $\mathcal{Q}_t$ enter the layer; the wallet masks candidates to $\mathcal{A}_t$ and exposes pressure $\lambda_t$; \padct{} selects the paid arm; the execution substrate returns a response and receipt; and the layer passes chosen-arm feedback $(\ut,c_t)$ to the policy.}
\Description{Architecture diagram showing task intake, quote filtering by wallet affordability, PA-DCT provider selection, payment execution, receipt collection, and policy feedback update.}
\label{fig:architecture-short}
\end{figure*}

402Pilot turns this formulation into a boundary between the autonomous agent and the payment substrate (Fig.~\ref{fig:architecture-short}). On the payment side, it consumes quoted candidates, masks infeasible actions, delegates the selected payment and service call, records spend, and receives a receipt. On the policy side, it exposes a small interface: \texttt{select}$(\mathbf{x}_t,\mathcal{A}_t)\to a_t$ and \texttt{update}$(\mathbf{x}_t,a_t,\ut,c_t)$. The layer does not modify x402, hold funds, or implement settlement. Its purpose is to turn payment execution into a learning interface: the substrate constrains what can be bought, and the policy learns what is worth buying.

\section{\texorpdfstring{The \padct{} Policy}{The PA-DCT Policy}}
\label{sec:method}

\padct{} instantiates the policy side of 402Pilot. For each provider $a$ and task bucket $k$, it maintains two Normal--Normal posteriors: a Q-posterior over utility and a C-posterior over realized cost. In the benchmark, $k(\mathbf{x}_t)$ maps tasks to four buckets: coding, multi-hop QA, closed-form web QA, and open-ended QA. Each cell stores an effective count $n^{a,k}$, a utility sum $S^{a,k}$, and a cost sum $S_c^{a,k}$. The Q-posterior is
\begin{equation}
\label{eq:posterior-short}
\begin{aligned}
\Sigmaq{a}{k}
&=\left(\frac{1}{\sigma^2_{0,q}}+\frac{n^{a,k}_t}{\sigma^2_q}\right)^{-1},\\
\thetaq{a}{k}
&=\Sigmaq{a}{k}\left(\frac{\mu_{0,q}}{\sigma^2_{0,q}}+\frac{S^{a,k}_t}{\sigma^2_q}\right),
\end{aligned}
\end{equation}
and the C-posterior uses the same form with $(\thetac{a}{k},\Sigmac{a}{k})$, prior mean $\bar c_a$, and cost sum $S_c^{a,k}$. The learned C-posterior is what lets \padct{} respond to price drift: if prices were fixed and visible before selection, the specification price would be sufficient; in the 402Pilot contract, effective charges are learned from receipts.

To track non-stationarity, \padct{} discounts every cell at the start of each round:
\begin{equation}
n^{a,k}_t=\gamma n^{a,k}_{t-1},\qquad
S^{a,k}_t=\gamma S^{a,k}_{t-1},\qquad
{S_c}^{a,k}_t=\gamma {S_c}^{a,k}_{t-1}.
\end{equation}
After the chosen call, only the selected cell receives the new observation. At selection time, \padct{} samples utility and cost from the corresponding posteriors for each affordable arm and scores
\begin{equation}
\label{eq:score-short}
s_a=(1-\lamn)\hatu{a}{k_t}
  - \lamn\,\mathrm{clip}(\hatc{a}{k_t}/c_{\max},0,1).
\end{equation}
The policy chooses the affordable arm with the largest sampled score. Wallet pressure shifts the decision toward cost control when burn rate is high, discounting lets stale evidence fade after market changes, and posterior sampling preserves exploration. All experiments use a single locked configuration: $\gamma=0.999$, $\nu=0.5$, $\alpha=2.0$, and $c_{\max}=\$0.01$.

\begin{algorithm}[t]
\caption{\padct{} policy loop.}
\label{alg:padct-short}
\footnotesize
\begin{algorithmic}[1]
\Require Spec costs $\bar c_a$; discount $\gamma$; normalizer $c_{\max}$.
\For{each served round $t$ with $\mathcal{A}_t\neq\emptyset$}
  \State Discount all cell counts and sums by $\gamma$
  \State $k_t\leftarrow k(\mathbf{x}_t)$;\quad $\lamn\leftarrow\lambda_t/(1+\lambda_t)$
  \For{$a\in\mathcal{A}_t$}
    \State Sample $\hatu{a}{k_t}$ and $\hatc{a}{k_t}$ from Q/C posteriors
    \State $s_a\leftarrow(1-\lamn)\hatu{a}{k_t}-\lamn\,\mathrm{clip}(\hatc{a}{k_t}/c_{\max},0,1)$
  \EndFor
  \State Select $a_t\leftarrow\arg\max_{a\in\mathcal{A}_t}s_a$
  \State Observe $(\ut,c_t)$ and update only cell $(a_t,k_t)$
\EndFor
\end{algorithmic}
\end{algorithm}

\section{Evaluation}
\label{sec:eval}

The evaluation asks whether a buyer-side decision layer helps when reliability or prices change after deployment. Live provider calls would add cost and confound policy effects with API drift, so 402Pilot-Bench uses frozen replay: responses are pre-generated and scored for each task--provider pair, and all policies replay the same task order, response-version draws, wallet budget, and scenario events under paired seeds.

\paragraph{Benchmark, providers, and scenarios.}
The benchmark contains 823 effective tasks across coding (HumanEval~\citep{chen2021humaneval}), multi-hop QA (HotpotQA~\citep{yang2018hotpotqa}), closed-form web QA (TriviaQA-web~\citep{joshi2017triviaqa}), and open-ended QA (OpenAssistant~\citep{kopf2023openassistant}). It stores five response versions for each task--provider pair, yielding 20,575 scored responses after one documented TriviaQA exclusion. The five providers form a price--quality ladder (P-cheap/P-mid/P-premium) plus a same-price reference/decoy trio: P-mid is reliable, P-adv produces fluent wrong answers, and P-flaky returns billed timeouts. A selector that only reads static metadata cannot separate these three arms.

The scenarios are S1 stationary operation, S2 a P-mid reliability outage from rounds 3,000--5,500, and S3 a P-premium price drop from \$0.01 to \$0.002 at round 1,000. Each policy--scenario cell uses 30 paired seeds, 10,000 rounds, and a \$50 budget; at the original premium price, this budget buys only 5,000 premium calls, so a policy must manage wallet pressure rather than simply choosing the strongest model throughout.

\begin{table}[t]
\centering
\caption{Provider set used in 402Pilot-Bench. P-mid, P-adv, and P-flaky share the same price tier and base model.}
\label{tab:providers-short}
\footnotesize
\setlength{\tabcolsep}{2pt}
\begin{tabularx}{\linewidth}{@{}lcX@{}}
\toprule
Provider & Price & Pipeline / role \\
\midrule
P-cheap   & $0.0005$ & qwen3.5-flash; uniform prompt; weakest cheap tier \\
P-mid     & $0.002$  & GPT-5.4-mini; uniform prompt; reliable mid tier \\
P-premium & $0.01$   & GPT-5.4; uniform prompt; strongest premium tier \\
P-adv     & $0.002$  & GPT-5.4-mini; 60\% adversarial prompt; fluent wrong answers \\
P-flaky   & $0.002$  & GPT-5.4-mini; 40\% timeout injection; billed timeouts \\
\bottomrule
\end{tabularx}
\end{table}

\paragraph{Comparators and metrics.}
We compare \padct{} against Random, fixed-arm policies, BudgetRule, Contextual DS-TS, Contextual BTS, and a replay True Oracle. BudgetRule is a static cascade-style heuristic: premium above 50\% remaining budget, mid above 20\%, and cheap below that. Contextual DS-TS tracks discounted utility but lacks wallet pressure and cost learning; Contextual BTS learns utility and cost but does not discount and scores by a raw utility/cost ratio. Metrics are mean quality $\bar q_T$, budget use, ROI, and PA-gap/$T$ to the replay True Oracle. ROI is an efficiency diagnostic rather than the objective: a cheap policy can have high ROI while serving low-quality answers. All means average over the same 30 seed identifiers within each scenario, and claimed differences use paired-seed tests.

\begin{table*}[t]
\centering
\small
\setlength{\tabcolsep}{3.4pt}
\caption{Main results on 402Pilot-Bench, averaged over 30 seeds. Columns report mean quality, budget used (\%), ROI, and PA-gap/$T$ to the replay True Oracle (lower is better). Bold marks the best non-oracle entry in rank-comparable quality, ROI, and PA-gap columns; budget use is descriptive.}
\label{tab:main-short}
\resizebox{\linewidth}{!}{%
\begin{tabular}{l|cccc|cccc|cccc}
\toprule
& \multicolumn{4}{c|}{S1 Stationary} & \multicolumn{4}{c|}{S2 Mid Outage} & \multicolumn{4}{c}{S3 Premium Promo} \\
Policy & $\bar q_T$ & budget & ROI & Gap & $\bar q_T$ & budget & ROI & Gap & $\bar q_T$ & budget & ROI & Gap \\
\midrule
Random            & $0.688$ & $66$ & $209$ & $0.365$ & $0.676$ & $66$ & $205$ & $0.375$ & $0.688$ & $37$ & $370$ & $0.274$ \\
Always-P-cheap    & $0.610$ & $10$ & $\mathbf{1220}$ & $0.167$ & $0.610$ & $10$ & $\mathbf{1220}$ & $\mathbf{0.164}$ & $0.610$ & $10$ & $\mathbf{1220}$ & $0.195$ \\
Always-P-mid      & $0.819$ & $40$ & $410$ & $\mathbf{0.101}$ & $0.757$ & $40$ & $379$ & $0.174$ & $0.819$ & $40$ & $410$ & $0.129$ \\
Always-P-premium  & $0.433$ & $100$ & $87$ & $1.072$ & $0.433$ & $100$ & $87$ & $1.070$ & $\mathbf{0.865}$ & $56$ & $309$ & $0.400$ \\
BudgetRule        & $\mathbf{0.831}$ & $80$ & $208$ & $0.692$ & $\mathbf{0.769}$ & $80$ & $192$ & $0.722$ & $0.859$ & $56$ & $307$ & $0.405$ \\
\midrule
Contextual DS-TS  & $0.559$ & $100$ & $112$ & $0.858$ & $0.514$ & $100$ & $103$ & $0.884$ & $0.840$ & $48$ & $349$ & $0.264$ \\
Contextual BTS    & $0.612$ & $10$ & $1187$ & $0.169$ & $0.612$ & $10$ & $1187$ & $0.166$ & $0.612$ & $10$ & $1187$ & $0.197$ \\
\midrule
\textbf{\padct{} (ours)} & $0.797$ & $42$ & $378$ & $0.133$ & $0.761$ & $43$ & $357$ & $0.166$ & $0.831$ & $39$ & $429$ & $\mathbf{0.121}$ \\
True Oracle       & $0.901$ & $32$ & $561$ & $0.000$ & $0.901$ & $33$ & $551$ & $0.000$ & $0.906$ & $25$ & $722$ & $0.000$ \\
\bottomrule
\end{tabular}%
}
\end{table*}

\begin{figure}[t]
\centering
\includegraphics[width=0.9\linewidth]{figures/S1.png}
\vspace{-0.5ex}
\includegraphics[width=0.9\linewidth]{figures/S2.png}
\vspace{-0.5ex}
\includegraphics[width=0.9\linewidth]{figures/S3.png}
\caption{\padct{} provider-allocation share over rounds. Shaded regions mark shock windows: the policy remains mid-dominated in S1, moves away from P-mid during S2, and shifts toward P-premium in S3.}
\Description{Three stacked line charts showing PA-DCT allocation shares across providers in stationary, mid-outage, and premium-promotion scenarios.}
\label{fig:learning-curves-short}
\end{figure}

\paragraph{Main results.}
Table~\ref{tab:main-short} shows that \padct{} is an adaptive trade-off rather than a one-column winner. In stationary S1, the reliable mid-tier arm is strong and \padct{} pays an exploration cost: Always-P-mid reaches the lowest non-oracle PA-gap/$T$ while \padct{} trails by about 2.2 quality points. In S2, the policy learns around the mid-tier outage and reaches $\bar q_T=0.761$ while spending 43\% of the wallet, compared with 0.610 for the cheapest fixed policy and 0.757 for Always-P-mid. BudgetRule slightly improves quality, but spends 80\% of the budget and suffers a large PA-gap because it follows fixed thresholds rather than feedback. In S3, \padct{} captures the premium promotion, improving over Always-P-mid on quality (0.831 vs.\ 0.819), ROI (429 vs.\ 410), and PA-gap/$T$ (0.121 vs.\ 0.129).

The learning baselines sharpen the failure modes. Contextual DS-TS tracks utility drift but lacks wallet pressure and exhausts the budget in S1/S2. Contextual BTS learns cost but does not discount old evidence and optimizes a raw utility/cost ratio, so under the 20$\times$ cost spread it locks onto P-cheap and misses the S3 opportunity. The allocation curves in Fig.~\ref{fig:learning-curves-short} show why the combined design matters: \padct{} moves away from P-mid during the outage, shifts toward P-premium after the price drop, and stays mid-dominated when no shock occurs.

\paragraph{Robustness across axes.}
Taken as a trade-off panel, Table~\ref{tab:main-short} shows that \padct{} is the only non-oracle policy without a clear failure mode across the three markets. Cheap policies and Contextual BTS conserve budget but under-serve tasks; premium-heavy policies and BudgetRule can buy quality in the wrong regime and pay a large PA-gap; Contextual DS-TS tracks quality drift but lacks wallet pressure and exhausts the budget in S1/S2. As a compact robustness descriptor, \padct{} has the lowest mean rank (2.89/8, lower is better) and the smallest worst rank (4/8) across the nine rank-comparable cells in Table~\ref{tab:main-short}. The per-axis numbers remain the load-bearing evidence, but this summary matches the central finding: the decision layer matters because the buyer must keep learning and adapting under market change, not merely win one static metric.

\paragraph{Quality under a budget.}
The learned reallocations matter because the cheap safe policy under-serves tasks, while static high-quality policies can spend in the wrong places. In S2, \padct{} gains 15.1 quality points over Always-P-cheap while spending only 43\% of the wallet, and its smaller gain over Always-P-mid is significant under paired seeds ($\Delta\bar q_T=+0.004$, $p=1.67\times10^{-3}$). In S3, \padct{} improves over Always-P-mid on quality, ROI, and PA-gap/$T$ at the same time. Always-P-premium reaches the highest S3 quality once premium becomes cheap, but it does so with lower ROI and much worse PA-gap because it pays for premium before the promotion as well.

\paragraph{Statistical checks.}
Paired-seed tests match these claims. The S3 improvements over Always-P-mid on quality, ROI, and PA-gap/$T$ are all significant ($p<10^{-20}$), and \padct{} has a significantly smaller PA-gap/$T$ than both learning baselines in S1 and S3. The S1 disadvantages against Always-P-mid and BudgetRule are also significant, which is why we treat them as the exploration cost of maintaining shock-ready posteriors rather than hide them. The one borderline cell is S2 PA-gap/$T$: \padct{} is numerically close to Always-P-cheap and Contextual BTS, while delivering much higher quality than those cheap-biased policies.

\begin{table}[t]
\centering
\caption{Component ablations. Entries are PA-gap/$T$; lower is better.}
\label{tab:ablation-mini}
\footnotesize
\setlength{\tabcolsep}{4pt}
\begin{tabularx}{\columnwidth}{@{}>{\raggedright\arraybackslash}Xccc@{}}
\toprule
Variant & S1 & S2 & S3 \\
\midrule
Full \padct{} & $0.133$ & $0.166$ & $0.121$ \\
$-P$ no pressure & $0.855$ & $0.878$ & $0.259$ \\
$-D$ no discount & $0.109$ & $0.170$ & $0.114$ \\
$-TS$ greedy & $0.130$ & $0.175$ & $0.145$ \\
$-C_{\mathrm{post}}$ spec cost & -- & -- & $0.144$ \\
\bottomrule
\end{tabularx}
\end{table}

\paragraph{Ablation summary.}
Table~\ref{tab:ablation-mini} shows which design choices fail under which market. The most severe failure is removing wallet pressure: the policy exhausts the wallet in S1/S2 and its PA-gap rises by more than sixfold. Discounting and cost learning matter in the shock regimes: without discounting, reliability recovery slows in S2, while a spec-only cost estimate misses the S3 promotion. Greedy selection is less stable across seeds and misses S3 in several runs; contextual bucketing mainly protects task-specific separation.

\section{Discussion and Conclusion}
\label{sec:discussion}

402Pilot does not replace payment protocols; it supplies the missing buyer policy above them. The boundary is deliberately narrow: it assumes posted payable candidates, chooses only from the affordable set, learns only from chosen outcomes, and leaves discovery, authorization, settlement faults, and observability to the surrounding payment stack. This makes the learning problem compatible with x402-style execution while making wallet pressure and receipt-revealed cost first-class decision signals.

The main limitation is that the quantitative evaluation is controlled. Frozen replay makes policy comparisons reproducible, but it does not measure live settlement latency, facilitator failures, congestion, or provider drift in deployment. The benchmark also covers a finite set of task types, providers, and scripted shocks.

Overall, the evidence supports a simple conclusion: autonomous payment systems need learning at the buyer layer, not only safer settlement mechanisms. \padct{} is not universally optimal in a static market, but it is the policy in this comparison that keeps service quality, spending efficiency, and payment-aware objective jointly competitive when reliability and prices change.

\paragraph{Code and supplement.}
The companion technical report contains the supplementary notation, proof scope, computation notes, detailed ablations, sensitivity sweeps, and local x402 quote--pay--receipt witness:
\url{https://github.com/MCCodeAI/402Pilot/raw/main/docs/technical_report.pdf}.
Code and evaluation materials are available at \url{https://github.com/MCCodeAI/402Pilot}.

\paragraph{Author contributions.}
Yin Li led ideation, method design, experiments, implementation, analysis, and writing. Yanbo He contributed to paper discussions, result audit, and baseline optimization. Wanting Xiong contributed related work, project discussion, and audit.

\bibliographystyle{ACM-Reference-Format}
\bibliography{references}

\end{document}
\endinput
